% Горизонты: Диалог наук — шаблон рукописи (LaTeX), v3
\documentclass[12pt,a4paper]{article}
\usepackage[T2A]{fontenc}
\usepackage[utf8]{inputenc}
\usepackage[russian,english]{babel}
\usepackage[margin=2.5cm]{geometry}
\usepackage{setspace}
\onehalfspacing

\title{Название статьи}
\author{И.\,О. Фамилия$^{1}$, И.\,О. Фамилия$^{2}$\\
{\small $^{1}$Организация, город; $^{2}$Организация, город}}
\date{}

\begin{document}
\maketitle

\begin{abstract}
Аннотация (150--250 слов): цель, методы, результаты, выводы.\\[6pt]
\textbf{Ключевые слова:} слово 1, слово 2, слово 3, слово 4, слово 5.
\end{abstract}

\section{Введение}
Текст.

\section{Методы}
Текст.

\section{Результаты}
Текст.

\section{Обсуждение}
Текст.

\begin{thebibliography}{99}
\bibitem{1} Список литературы оформляется по ГОСТ Р 7.0.100--2018 или APA 7.
\end{thebibliography}
\end{document}
